\subsection{TOPOLOGY OF VECTOR SPACES AND ALGEBRAS OVER THE FIELD R}

Let \(\mathbf{E}\) be a vector space of dimension \(n\) over the field \(\mathbf{R}\) ; if \((\boldsymbol{a}_i)_{1 \leqslant i \leqslant n}\) is a basis of \(\mathbf{E}\) , then every point \(\boldsymbol{x} \in \mathbf{E}\) can be written uniquely in the form \(\boldsymbol{x} = \sum_{i=1}^{n} x_i \boldsymbol{a}_i\) , where the \(x_i\) are real numbers; the mapping \((x_i) \to \sum_{i=1}^{n} x_i \boldsymbol{a}_i\) is therefore a bijective linear mapping of \(\mathbf{R}^n\) onto \(\mathbf{E}\) . If we transport to \(\mathbf{E}\) the topology of \(\mathbf{R}^n\) by this mapping, \(\mathbf{E}\) is endowed with a topology compatible with its additive group structure, and the mapping \((t, \boldsymbol{x}) \to tx\) of \(\mathbf{R} \times \mathbf{E}\) into \(\mathbf{E}\) is continuous with respect to this topology. This topology is independent of the basis chosen in \(\mathbf{E}\) ; for if \((\boldsymbol{a}')\) is another basis of \(\mathbf{E}\) , and if \(x = \sum_{i=1}^{n} x_i' \boldsymbol{a}_i' = \sum_{i=1}^{n} x_i \boldsymbol{a}_i\) , the mapping \((x_i) \to (x_i')\) of \(\mathbf{R}^n\) onto itself is linear and therefore a homeomorphism.

This fact leads one to suspect that the topology so defined on E should be capable of characterization without the help of a basis of E. In fact, we shall see later that this is the only Hausdorff topology on E for which the functions x - y (on E × E) and tx (on R × E) are continuous.

If now A is an algebra of finite rank n over the field R, the above topology on A (considered as an n-dimensional vector space over R) is compatible not only with the additive group structure of A, but also with its ring structure. This is a consequence of the following more general result:

PROPOSITION 5. Let E, F, G be three finite-dimensional vector spaces over the field R. Then every bilinear mapping (*) f of E × F into G is continuous.

We may suppose that \(E = R^{m}\) , \(F = R^{n}\) , \(G = R^{p}\) ; it is enough to show that the coordinates in \(R^{p}\) of \(f(x, y)\) are continuous functions of \((x, y) \in E \times F\) (Chapter I, § 4, no. 1, Proposition 1). In other words, it is enough to show that every bilinear form g is continuous on \(E \times F\) ; and this is immediate, since \(g(x, y)\) is a polynomial in the coordinates of x and y.

